\begin{description}
\item[a.] The light curve. \hfill(10 points)
  % FIGURE NEEDED: plotted light curve of HV 2257, V (mag) vs phase, from
  % Gieren 1993 (MNRAS vol 265); 2015 DA solutions, p. 1/9.

\item[b.] The color curve. \hfill(10 points)
  % FIGURE NEEDED: plotted colour curve of HV 2257, (V-R)_J vs phase, from
  % Gieren 1993 (MNRAS vol 265); 2015 DA solutions, p. 1/9.

\item[c.] The radial velocity curve. \hfill(10 points)
  % FIGURE NEEDED: plotted radial-velocity curve of HV 2257, radial velocity
  % (km/s) vs phase, from Gieren 1993 (MNRAS vol 265); 2015 DA solutions,
  % p. 2/9.
\end{description}

\noindent\textbf{Solar absolute bolometric magnitude calculation}
\hfill(5 points)

$F_\odot$ is the solar flux at the distance 10 pc which corresponds to the
solar bolomatric absolute magnitude.
\[ F_\odot = \frac{L_\odot}{4\pi \cdot (10\,pc)^2}
   = \frac{3.96 \cdot 10^{26}}
          {4 \cdot 3.14 \cdot 100 \cdot 3.086^2 \cdot 10^{32}}
   = 3.31 \cdot 10^{-10}\,\mathrm{W/m^2} \]

\medskip
\noindent\textbf{Formula derivation} \hfill(15 points)

From Stefan--Boltzmann equation we have luminosity of star:
\[ L = 4\pi R^2 \sigma T^4, \]
here $R$ is radius of the star, $\sigma$ is the Stefan--Boltzmann constant,
and $T$ is effective temperature of star. Star flux at distance $d$ will be
equal:
\[ F = \frac{\sigma R^2 T^4}{d^2} \]
From Pogson definition we have:
\[ m_1 - m_2 = -2.5\log\frac{F_1}{F_2} \]
Or using the Sun as reference, a star's observed bolometric flux will be:
\[ F = F_\odot\,10^{-\frac{m_{bol} - M_{\odot bol}}{2.5}} \]

Consider two moments, say $t_1$ and $t_2$. It is better to choose the phase
$t_1$ and $t_2$ during which the star's expansion acceleration close to
constant, and the difference in magnitude and color as large as possible. At
the moment $t_1$, with measured temperature $T_1$ and radius $R_1$, absolute
bolometric flux will be:
\[ F_1 = \frac{\sigma R_1^2 T_1^4}{d^2}
   \tag*{(1a)} \]
Later, at moment $t_2$:
\[ F_2 = \frac{\sigma R_2^2 T_2^4}{d^2}
   \tag*{(1b)} \]
During this time, the star's atmosphere has expanded from $R_1$ to $R_2$:
\[ R_2 = R_1 + \Delta R \]
or:
\[ \frac{R_2}{R_1} = 1 + \frac{\Delta R}{R_1} \]
Reminding:
\[ \frac{F_2}{F_1}
   = \left( \frac{T_2}{T_1} \right)^4 \left( \frac{R_2}{R_1} \right)^2 \]
We have
\[ \frac{F_2}{F_1}
   = \left( \frac{T_2}{T_1} \right)^4
     \left( \frac{\Delta R}{R_1} + 1 \right)^2
   \tag*{(2)} \]
From Pogson definition, and (2) we can find:
\[ R_1 = \frac{\Delta R}
   {\left( \frac{T_1}{T_2} \right)^2 \cdot 10^{-\frac{m_2 - m_1}{5}} - 1} \]

\medskip
\noindent\textbf{Calculation of $\Delta R$ from the radial velocity graph}
\hfill(20 points)

For finding $\Delta R$, we can use radial velocity curve, taking two moments,
$t_1$ and $t_2$, between which the expansion acceleration can be assumed
constant,
\[ \Delta R = \frac{(v_2 - v_1)\,\Delta\varphi\,P}{2} \]
here $P$ is pulsation period of star and $\Delta\varphi$ is phase difference
in moment between $t_1$ and $t_2$.

\begin{quote}
Alternative method using integral / sum:
\[ \Delta R = \int_{t_1}^{t_2} v(t)\,dt \]
We can calculate the integral by drawing lines connecting two adjacent points,
calculating the area of the trapesium under the line segment and sum up for
all line segment between moment $t_1$ and $t_2$.
\end{quote}

So, from radial velocity curve we choose the part of the graph which is close
to linear:
% FIGURE NEEDED: radial-velocity curve with the chosen near-linear segment
% between phases 0.75 and 0.85 marked; 2015 DA solutions, p. 4/9.
\begin{align*}
  t_1 &= 0.75 \\
  t_2 &= 0.85
\end{align*}
Which corresponds to
\begin{align*}
  v(t_1) &= 272\ \mathrm{km/s} \\
  v(t_2) &= 230\ \mathrm{km/s}
\end{align*}
Now we can calculate $\Delta R$:
\[ \Delta R = \frac{(230000 - 272000)(0.85 - 0.75) \cdot 39.294 \cdot 86400}
   {2} = -7.13 \cdot 10^9\,\mathrm{m} \]

\medskip
\noindent\textbf{Temperatures and magnitudes from photometric data}
\hfill(10 points)

For moments $t_1$ and $t_2$: from light curve we have the magnitude $V$:
\begin{align*}
  V_1 &= 13.58 \\
  V_2 &= 12.55
\end{align*}
From color curve we have the color $V-R$:
\begin{align*}
  (V-R)_1 &= 0.81 \\
  (V-R)_2 &= 0.53
\end{align*}
From Fig.~1 we can find the temperatures:
\begin{align*}
  T_1 &= 4000\,\mathrm{K} \\
  T_2 &= 4750\,\mathrm{K}
\end{align*}
From table 4 we have bolometric corrections for these moments (by using
Table 4 with linear interpolation):
\begin{align*}
  BC_1 &= -1.23 \\
  BC_2 &= -0.52
\end{align*}
Now we can calculate bolometric magnitudes:
\begin{align*}
  m_{bol}(t_1) &= 13.58 - 1.23 = 12.35 \\
  m_{bol}(t_2) &= 12.55 - 0.52 = 12.03
\end{align*}

\medskip
\noindent\textbf{Calculate $R_1$} \hfill(10 points)

First we calculate star's observed flux:
\[ F_1 = 2.913 \cdot 10^{-10} \cdot 10^{-\frac{12.35 - 4.72}{2.5}}
   = 2.584 \cdot 10^{-13}\,\mathrm{W/m^2} \]
% sic — the source uses F_sun = 2.913e-10 here although the value computed
% above was 3.31e-10 W/m^2.
Then the radius of the star at the moment $t_1$ will be:
\[ R_1 = \frac{\Delta R}
   {\left( \frac{T_1}{T_2} \right)^2 \cdot 10^{-\frac{m_2 - m_1}{5}} - 1}
   = \frac{-7.13 \cdot 10^9}
   {\left( \frac{4000}{4750} \right)^2
    \cdot 10^{-\frac{12.03 - 12.35}{5}} - 1}
   = 4.0 \cdot 10^{10}\,\mathrm{m} \]

\medskip
\noindent\textbf{Calculate distance} \hfill(10 points)

\[ d = \sqrt{ \frac{\sigma \cdot R_1^2 \cdot T_1^4}{F_1} }
   = 3.0 \cdot 10^{20}\,\mathrm{m}
   = \frac{3.0 \cdot 10^{20}}{3.086 \cdot 10^{16}} = 9.72\ \mathrm{kpc} \]

Answer: 9.72 kpc
